\chapter{Pruebas de hipótesis}
\labch{pruebas-hipotesis}

\begin{ideanote}
La muestra, ¿nos permite afirmar nuestras ideas previas? ¿No nos permite? ¿Cuánto tiene que ser la diferencia entre lo que sospechábamos y lo que nos dice la muestra para decir que no es lo mismo?
\end{ideanote}

Una de las razones por las que la estadística es útil es porque ayuda en la vida cotidiana o laboral. Una de las formas en que es útil es porque nos ayuda a tomar decisiones. Una manera de tomar decisiones es mediante pruebas de hipótesis.

Una hipótesis es una afirmación sobre la población. En nuestro caso, haremos pruebas de hipótesis \emph{paramétricas,} es decir, afirmaciones sobre algún parámetro de la población.

Una \emph{prueba de hipótesis} es una manera sistemática de comprobar si una hipótesis puede ser plausible o no. Esta comprobación estará basada en nuestros datos, es decir, en la muestra que tenemos. Por supuesto, esto involucra azar y por lo tanto podría ocurrir que tomemos una muestra que, mediante la intervención de los dioses del azar, nos dé un estimador del parámetro que en realidad no se parece al de la población. Sin embargo, podemos hacer algunos cálculos para medir la probabilidad de que eso ocurra.

Las pruebas que realizaremos seguirán un proceso clásico:
\begin{enumerate}
\item Definiremos dos hipótesis que contrastaremos, la hipótesis nula y la alternativa.
\item Definiremos un nivel de significancia.
\item Calcularemos un estadístico de prueba.
\item \begin{itemize}
\item O bien, a partir del nivel de significancia definiremos una región de rechazo, y veremos si el estadístico de prueba cae en ella.
\item O bien calcularemos un valor p a partir del estadístico de prueba.
\end{itemize}
\item Basado en el paso anterior, tomaremos una decisión acerca de nuestra hipótesis nula.
\end{enumerate}

\section{Definición de hipótesis}
\labsec{definicion-hipotesis}

Hay algunas sutilezas que tener en cuenta. Por ejemplo, supongamos que estamos haciendo una medición sobre la cantidad de bebida en las botellas y queremos comparar si tienen lo que dicen tener, \qty{500}{\milli\liter}. 

La hipótesis nula se refiere generalmente a la hipótesis que mantiene todo tal cual, en elegante, la hipótesis \emph{más parsimoniosa.} En nuestro caso, podríamos decir que la hipótesis nula es efectivamente que las latas contienen, en promedio, \qty{500}{\milli\liter}.

La hipótesis nula es lo único que podemos decir realmente, lo que significa que podemos aceptarla o bien, rechazarla. ¿En qué sentido es \enquote{lo único que podemos decir realmente}? En el contexto en que estamos trabajando, la hipótesis alternativa se constituye como una \emph{negación} de la hipótesis nula.\sidenote{y no es, por ejemplo, una hipótesis como decir \(\mu = 490\).} Esta negación la entenderemos en un sentido un poco más amplio que en la lógica aristotélica: si nuestra hipótesis nula es que el promedio de la población es \qty{500}{\milli\liter}, eso es decir que \[H_0: \mu = 500\] y nuestra hipótesis alternativa será una de las siguientes \[H_a: \mu<500,\quad H_a: \mu \neq 500,\quad H_a: \mu>500.\]

¿Cuál hipótesis alternativa es la que usaremos? Eso es parte de nuestra decisión y depende de nuestra intención. Si queremos demostrar que la compañía embotelladora está mintiendo (y ganar una jugosa demanda), probablemente usaríamos la primera versión, ya que dice que la verdadera media de la población es menor que lo que dicen las latas. Si somos un poco más neutrales, la segunda versión está bien, no nos interesa demostrar que es \emph{menor,} solo que no es \qty{500}{\milli\liter} (dejando abierta la posibilidad que sea mayor). Si somos los abogados de la embotelladora, nos interesaría usar la tercera y demostrar que la compañía da más bebida de la que promete.

\formal
Técnicamente, las hipótesis nula y alternativa \emph{deben} ser complementarias, es decir, si mi hipótesis alternativa es \(H_a: \mu>500\), mi hipótesis nula \emph{debe} ser \(H_0: \mu \leq 500\). Sin embargo, podemos ser no tan estrictos.

\informal
\section{Significancia, potencia y otras hierbas}
\labsec{significancia}

El segundo paso dice \enquote{Definiremos un nivel de significancia}, ¿pero qué significa eso?

Para empezar, necesitamos reconocer algo: nunca sabremos si \emph{de verdad} la hipótesis nula es verdadera o falsa. Para realmente saberlo, necesitaríamos conocer toda la población y ya sabemos que eso es generalmente imposible. Lo mejor que podemos hacer es decir que nuestros datos \emph{nos sugieren que} la hipótesis nula es verdadera o falsa.

Pero eso nos plantea una situación como en la \reftab{hachecero}. La hipótesis nula puede ser verdadera o falsa e indepedientemente de esto, basados en nuestros datos, nosotros podemos aceptarla o rechazarla.

\begin{table}[ht]
  \centering
  \begin{tabular}{>{\raggedright\arraybackslash}m{2.5cm}>{\centering\arraybackslash}m{1.5cm}>{\centering\arraybackslash}m{1.5cm}}
    & \(H_0\) es verdadera & \(H_0\) es falsa \\
    \midrule
    Aceptamos \(H_0\) & Decisión correcta         & Error tipo II (\(\beta\)) \\\midrule
    Rechazamos \(H_0\) & Error tipo I (\(\alpha\)) & Decisión correcta \\
  \end{tabular}
  \caption{La verdad de \(H_0\) y nuestras decisiones.}
  \labtab{hachecero}
\end{table}

Si en la realidad \(H_0\) es verdadera y la aceptamos, o bien si en la realidad \(H_0\) es falsa y la rechazamos, tomamos la decisión correcta en cualquiera de los dos casos. Sin embargo, también podemos equivocarnos y podemos equivocarnos de dos formas:
\begin{itemize}
\item \(H_0\) puede ser cierta en la realidad y nosotros tomamos la decisión de rechazarla.
\item \(H_0\) puede ser falsa en la realidad y nosotros tomamos la decisión de no rechazarla.
\end{itemize}

Si rechazamos \(H_0\) cuando \(H_0\) es verdadera, llamamos a este tipo de error \emph{error tipo I.} El máximo de la probabilidad de cometer este error\sidenote{que podemos entender como un falso positivo} se define como \(\alpha\), la \emph{significancia} del test.

\formal
La definición de \(\alpha\) es \[\alpha = \sup_{\theta \in \Theta_0} \prob_{\theta}(T\in R),\] es decir, \(\alpha\) es el supremo sobre \(\theta\) de la probabilidad de que el estadístico de prueba pertenezca a la región de rechazo.

\informal
Por otra parte, el error que cometemos al aceptar \(H_0\) cuando en la realidad es falsa es llamado el \emph{error tipo II.} Definiremos su probabilidad como \(\beta\). 

¿Cómo se ven \(\alpha\) y \(\beta\)? Por ejemplo, si nuestra hipótesis alternativa es \(\mu>\mu_0\).


\begin{tikzpicture}
    \begin{axis}[
        width=1.3\linewidth,
        height=5cm,
        axis lines=middle,
        domain=-3.8:3.8,
        samples=150,
        ymin=0,
        enlargelimits=false,
        clip=false,
        ytick=\empty,
        xtick=\empty,
    ]
    \addplot[fill=DiceRed!30, draw=none, domain=1.645:3.8] {gauss(x,0,1)} \closedcycle;
    \addplot[DiceBlue!80!black, thick] {gauss(x,0,1)};
    \draw[gray] (axis cs:1.645, 0) -- (axis cs:1.645, {gauss(1.645,0,1)});
    \node[below, font=\small] at (axis cs:0, 0) {\(0\)};
    \node[below, font=\small] at (axis cs:1.645, 0) {\(z_\alpha\)};
    \node[DiceRed!50!black, font=\small] at (axis cs:2, 0.023) {\(\alpha\)};
    \end{axis}
\end{tikzpicture}

El error tipo 1 asume que \(H_0\) es verdadera. Por lo tanto, si estandarizamos como \[Z=\frac{\mean{X} - \mu_0}{\sigma/\sqrt{n}}\] podemos entender a \(\alpha\) como la probabilidad de que nos equivoquemos, es decir, que el promedio de nuestra muestra esté tan distante del (supuesto) verdadero promedio que sea inevitable pensar que la hipótesis nula es falsa.

Eso quiere decir que si la media muestral está lo suficientemente alejada de la media poblacional propuesta, tendremos razones para creer que \(H_0\) es falsa. ¿Qué significa \enquote{suficientemente alejada}? Significa que estamos dispuestos a equivocarnos: eso es lo que quiere decir la significancia.

¿Y cómo se ve \(\beta\)?

\begin{pausa}
Para calcular \(\alpha\), asumimos que \(H_0\) es verdadera y calculamos la probabilidad de caer en la región de rechazo. Para calcular \(\beta\), ¿qué necesitamos asumir además?
\end{pausa}

Para calcular \(\beta\), necesitamos especificar el verdadero valor del parámetro, es decir, un \(\mu_a\) concreto. Esto es distinto de \(\alpha\), cuyo cálculo no requiere conocer la alternativa.

\formal
\[\beta(\mu_a) = \prob_{\mu_a}(T \notin R),\]
es decir, la probabilidad de que el estadístico de prueba \emph{no} caiga en la región de rechazo cuando el verdadero parámetro es \(\mu_a\).
\informal

Supongamos que el verdadero promedio es \(\mu_a > \mu_0\). La curva azul muestra cómo se distribuye el estadístico bajo \(H_0\) (centrada en \(\mu_0\)); la curva roja muestra la distribución real (centrada en \(\mu_a\)).

\begin{tikzpicture}
    \begin{axis}[
        width=1.3\linewidth,
        height=5cm,
        axis lines=middle,
        domain=-3.5:5.5,
        samples=150,
        ymin=0,
        enlargelimits=false,
        clip=false,
        ytick=\empty,
        axis y line = none,
        xtick=\empty,
    ]
    \addplot[fill=DiceBlue!30, draw=none, domain=-3.5:1.645] {gauss(x,2.5,1)} \closedcycle;
    \addplot[fill=DiceRed!30, draw=none, domain=1.645:5.5] {gauss(x,0,1)} \closedcycle;
    \addplot[DiceBlue!80!black, thick] {gauss(x,0,1)};
    \addplot[DiceRed!80!black, thick] {gauss(x,2.5,1)};
    \draw[gray] (axis cs:1.645, 0) -- (axis cs:1.645, {gauss(1.645,2.5,1)});
    \draw[gray!60, thin] (axis cs:0, 0) -- (axis cs:0, {gauss(0,0,1)});
    \draw[gray!60, thin] (axis cs:2.5, 0) -- (axis cs:2.5, {gauss(2.5,2.5,1)});
    \node[below, font=\small] at (axis cs:0, 0) {\(\mu_0\)};
    \node[below, font=\small] at (axis cs:2.5, 0) {\(\mu_a\)};
    \node[below, font=\small] at (axis cs:1.645, 0) {\(z_\alpha\)};
    \node[DiceRed!50!black, font=\small] at (axis cs:2, 0.023) {\(\alpha\)};
    \node[DiceBlue!50!black, font=\small] at (axis cs:0.9, 0.06) {\(\beta\)};
    \end{axis}
\end{tikzpicture}

\(\beta\) es la probabilidad de que, aunque la distribución real esté centrada en \(\mu_a\), el estadístico de prueba caiga a la izquierda de \(z_\alpha\) y \emph{no} rechacemos \(H_0\). Es la probabilidad de aceptar \(H_0\) aunque sea falsa.

Nótese que \(\beta\) depende de cuánto se aleja \(\mu_a\) de \(\mu_0\): si la diferencia es grande, las distribuciones se separan y es fácil detectar que \(H_0\) es falsa (\(\beta\) es pequeño). Si es pequeña, se superponen y es difícil distinguirlas (\(\beta\) es grande).

El complemento de \(\beta\) es la \emph{potencia} de la prueba, definida como \(1-\beta\). Una prueba con alta potencia detecta rápidamente cuando \(H_0\) es falsa.

\section{Estadístico de prueba}
\labsec{estadistico-prueba}

Nuestra muestra nos da datos. Dependiendo del parámetro que estemos usando en nuestra prueba de hipótesis, podremos calcular un estimador puntual para ese parámetro a partir de los datos de la muestra. 

El cálculo del estimador dependerá de la prueba que estemos realizando, de las maneras que ya conocemos. 
\begin{center}
{
\renewcommand{\arraystretch}{3}
\begin{tabular}{rl}
Prueba & Estadístico\\\midrule
\(\mu\), \(\sigma\) conocido & \(Z = \dfrac{\mean{x} - \mu}{\sigma/\sqrt{n}}\)\\
\(\mu\), \(\sigma\) desconocido & \(T = \dfrac{\mean{x} - \mu}{s/\sqrt{n}}\)\\
\(p\) & \(Z = \dfrac{\estim{p} - p}{\sqrt{\frac{p(1-p)}{n}}}\)\\
\(\sigma^2\) & \(\chi^2 = \dfrac{(n-1)s^2}{\sigma^2}\)
\end{tabular}
}
\end{center}

El estadístico de prueba es una representación de nuestra muestra en términos de la distribución (normal, t, chi cuadrado) que corresponde. La pregunta que queda por responder es
\begin{quotation}
Suponiendo que esta muestra es representativa, ¿el estadístico de prueba me permite creer que la hipótesis nula es razonable o no?
\end{quotation}

\section{Comparando al estadístico de prueba}
\labsec{comparando-estadistico-prueba}

Tenemos dos métodos para comparar. El clásico es hacerlo mediante región de rechazo. Sin embargo, usando un computador es extremadamente fácil hacerlo mediante el valor \(p\) (o \(p\)-valor).

\subsection{Región de rechazo}

Nuestro valor de \(\alpha\) es, recordemos, la probabilidad de que la muestra nos lleve a rechazar la hipótesis nula \emph{dado} que es verdadera. 

Si centramos, por ejemplo estandarizando una normal, a partir de la media \(\mu\) que nos da la hipótesis nula (como en el cálculo del estadístico de prueba para \(\mu\) con \(\sigma\) conocido), tenemos una situación así (para hipótesis alternativa de una sola cola).

\begin{tikzpicture}
    \begin{axis}[
        width=1.3\linewidth,
        height=5cm,
        axis lines=middle,
        domain=-3.8:3.8,
        samples=150,
        ymin=0,
        enlargelimits=false,
        clip=false,
        ytick=\empty,
        xtick=\empty,
    ]
    \addplot[fill=DiceRed!30, draw=none, domain=1.645:3.8] {gauss(x,0,1)} \closedcycle;
    \addplot[DiceBlue!80!black, thick] {gauss(x,0,1)};
    \draw[gray] (axis cs:1.645, 0) -- (axis cs:1.645, {gauss(1.645,0,1)});
    \node[below, font=\small] at (axis cs:0, 0) {\(0\)};
    \node[below, font=\small] at (axis cs:1.645, 0) {\(z_\alpha\)};
    \node[DiceRed!50!black, font=\small] at (axis cs:2, 0.023) {\(\alpha\)};
    \end{axis}
\end{tikzpicture}

Tenemos un cierto percentil \(z_{\alpha}\), que define dos regiones. Definiremos la \emph{región de rechazo} como aquella que nos lleva a rechazar la hipótesis nula si el estadístico de prueba cae dentro de ella. Como estamos dispuestos a rechazar la hipótesis nula, dado que ella es cierta, con probabilidad \(\alpha\), resulta que en el gráfico anterior, donde mostramos el mundo según la hipótesis nula, la región sombreada corresponde a la región de rechazo.

La \emph{región de no rechazo} es aquella que no nos lleva a rechazar la hipótesis nula. Por lo tanto, en el gráfico anterior es el área que no está marcada.

Debemos tener cuidado cuando se trate de una prueba de dos colas. En ese caso tenemos que considerar que nuestra probabilidad \(\alpha\) se refiere a que la verdadera media es \emph{distinta} de la propuesta en la hipótesis nula, lo que significa que puede ser mayor o menor. Al ser así, la región de rechazo queda especificada en \emph{dos} intervalos. Pero \(\alpha\) sigue siendo \(\alpha\), por lo que cada intervalo encierra una probabilidad de \(\alpha/2\).

\begin{tikzpicture}
    \begin{axis}[
        width=1.3\linewidth,
        height=5cm,
        axis lines=middle,
        domain=-3.8:3.8,
        samples=150,
        ymin=0,
        enlargelimits=false,
        clip=false,
        ytick=\empty,
        xtick=\empty,
    ]
    \addplot[fill=DiceRed!30, draw=none, domain=-3.8:-1.96] {gauss(x,0,1)} \closedcycle;
    \addplot[fill=DiceRed!30, draw=none, domain=1.96:3.8] {gauss(x,0,1)} \closedcycle;
    \addplot[DiceBlue!80!black, thick] {gauss(x,0,1)};
    \draw[gray] (axis cs:-1.96, 0) -- (axis cs:-1.96, {gauss(-1.96,0,1)});
    \draw[gray] (axis cs:1.96, 0) -- (axis cs:1.96, {gauss(1.96,0,1)});
    \node[below, font=\small] at (axis cs:0, 0) {\(0\)};
    \node[below, font=\small] at (axis cs:-1.96, 0) {\(-z_{\alpha/2}\)};
    \node[below, font=\small] at (axis cs:1.96, 0) {\(z_{\alpha/2}\)};
    \node[DiceRed!50!black, font=\small] at (axis cs:-2.6, 0.05) {\(\alpha/2\)};
    \node[DiceRed!50!black, font=\small] at (axis cs:2.6, 0.05) {\(\alpha/2\)};
    \end{axis}
\end{tikzpicture}

¿De qué lado queda la región de rechazo? Depende de la hipótesis alternativa, y conviene pensarlo en vez de memorizarlo. Si sospechamos que la verdadera media es \emph{mayor} que la propuesta (\(H_a: \mu>\mu_0\)), sólo nos van a hacer dudar de \(H_0\) las muestras con un promedio muy alto: la región de rechazo queda en la cola derecha. Si sospechamos que es \emph{menor} (\(H_a:\mu<\mu_0\)), queda en la cola izquierda. Y si sólo decimos que es \emph{distinta} (\(H_a:\mu\neq\mu_0\)), nos hacen dudar tanto los promedios muy altos como los muy bajos, así que la región se reparte en las dos colas, como acabamos de ver.

La regla, entonces, es simple: calculamos el estadístico de prueba con nuestra muestra y miramos dónde cae. Si cae en la región de rechazo, rechazamos \(H_0\); si cae en la región de no rechazo, no la rechazamos.

Volvamos a las botellas. Sospechamos que la embotelladora está rellenando de menos, así que contrastamos \(H_0: \mu=\qty{500}{\milli\liter}\) contra \(H_a: \mu<\qty{500}{\milli\liter}\), una prueba de una sola cola (la izquierda). Fijamos una significancia de \(\alpha=\num{0,05}\) y supongamos\sidenote{en la práctica casi nunca lo conocemos; cuando no, usamos el estadístico \(T\) con \(s\) en vez de \(\sigma\), pero la mecánica es idéntica} que conocemos la desviación estándar del proceso, \(\sigma=\qty{6}{\milli\liter}\). Medimos \(n=36\) botellas y obtenemos un promedio de \(\mean{x}=\qty{498}{\milli\liter}\).

El estadístico de prueba es \[Z = \frac{\mean{x}-\mu_0}{\sigma/\sqrt{n}} = \frac{498-500}{6/\sqrt{36}} = \frac{-2}{1} = -2.\]

Como la cola es izquierda y \(\alpha=\num{0,05}\), el percentil crítico es \(z_\alpha=-\num{1,645}\)\marginnote{lo sacamos de la tabla de la normal o del computador}. La región de rechazo es todo lo que esté a la izquierda de ese valor.

\begin{tikzpicture}
    \begin{axis}[
        width=1.3\linewidth,
        height=5cm,
        axis lines=middle,
        domain=-3.8:3.8,
        samples=150,
        ymin=0,
        enlargelimits=false,
        clip=false,
        ytick=\empty,
        xtick=\empty,
    ]
    \addplot[fill=DiceRed!30, draw=none, domain=-3.8:-1.645] {gauss(x,0,1)} \closedcycle;
    \addplot[DiceBlue!80!black, thick] {gauss(x,0,1)};
    \draw[gray] (axis cs:-1.645, 0) -- (axis cs:-1.645, {gauss(-1.645,0,1)});
    \draw[Mahogany, thick] (axis cs:-2, 0) -- (axis cs:-2, {gauss(-2,0,1)});
    \node[below, font=\small] at (axis cs:0, 0) {\(0\)};
    \node[below, font=\small] at (axis cs:-1.645, 0) {\(z_\alpha\)};
    \node[below, Mahogany, font=\small, yshift=-11pt] at (axis cs:-2, 0) {\(z^*\)};
    \node[DiceRed!50!black, font=\small] at (axis cs:-2.7, 0.04) {\(\alpha\)};
    \end{axis}
\end{tikzpicture}

Nuestro estadístico, \(z^*=-2\), es más pequeño que \(-\num{1,645}\): cae dentro de la región de rechazo. Por lo tanto, rechazamos la hipótesis nula. Con estos datos, y a un \qty{5}{\percent} de significancia, tenemos evidencia de que las botellas traen, en promedio, menos de \qty{500}{\milli\liter}.

\subsection{Valor \texorpdfstring{\(p\)}{p}}
\labsec{valor-p}

El método de la región de rechazo nos da un veredicto de sí o no: el estadístico cae adentro o no cae. Pero se guarda algo de información. No es lo mismo que el estadístico caiga justo pasando la frontera a que caiga lejísimos, en la punta de la cola. En ambos casos rechazamos, pero en el segundo deberíamos sentirnos mucho más seguros. El valor \(p\) es una forma de medir \emph{qué tan adentro} caímos.

\begin{pausa}
Si nuestro estadístico de prueba cayera \emph{justo} sobre la frontera de la región de rechazo, ¿qué valor cree que tomaría el valor \(p\)?
\end{pausa}

La idea es la siguiente. Nos paramos en el mundo de la hipótesis nula —es decir, suponemos que \(H_0\) es verdadera— y nos preguntamos: si \(H_0\) fuera cierta, ¿qué tan raro sería obtener una muestra como la nuestra, o todavía más extrema? Esa probabilidad es el valor \(p\).

En el ejemplo de las botellas, el estadístico fue \({z^*=-2}\). Como la prueba es de cola izquierda, \enquote{tan o más extremo} significa \enquote{tan o más a la izquierda}, es decir, un valor menor o igual a \(-2\). El valor \(p\) es el área bajo la curva a la izquierda de \(-2\):
\[p = \prob(Z \leq -2) = \num{0,0228}.\]

\begin{tikzpicture}
    \begin{axis}[
        width=1.3\linewidth,
        height=5cm,
        axis lines=middle,
        domain=-3.8:3.8,
        samples=150,
        ymin=0,
        enlargelimits=false,
        clip=false,
        ytick=\empty,
        xtick=\empty,
    ]
    \addplot[fill=DiceRed!30, draw=none, domain=-3.8:-2] {gauss(x,0,1)} \closedcycle;
    \addplot[DiceBlue!80!black, thick] {gauss(x,0,1)};
    \draw[Mahogany, thick] (axis cs:-2, 0) -- (axis cs:-2, {gauss(-2,0,1)});
    \node[below, font=\small] at (axis cs:0, 0) {\(0\)};
    \node[below, Mahogany, font=\small] at (axis cs:-2, 0) {\(z^*\)};
    \node[DiceRed!50!black, font=\small] at (axis cs:-2.7, 0.035) {\(p\)};
    \end{axis}
\end{tikzpicture}

En palabras: si la embotelladora estuviera diciendo la verdad y la media real fuera \qty{500}{\milli\liter}, sólo un \qty{2,28}{\percent} de las muestras de \num{36} botellas darían un promedio tan bajo (o más bajo) que el nuestro. O nos tocó una casualidad bastante grande, o la media real no es \num{500}.

\formal
El valor \(p\) es la probabilidad, calculada suponiendo que \(H_0\) es verdadera, de obtener un estadístico de prueba al menos tan extremo como el observado, \(z^*\). Qué significa \enquote{tan extremo} depende de la hipótesis alternativa:
\[p = \prob_{H_0}(Z \leq z^*) \quad(\text{cola izquierda}),\]
\[p = \prob_{H_0}(Z \geq z^*) \quad(\text{cola derecha}),\]
\[p = \prob_{H_0}(|Z| \geq |z^*|) = 2\,\prob_{H_0}(Z \geq |z^*|) \quad(\text{dos colas}).\]

\informal
En la prueba de dos colas un valor extremo hacia cualquiera de los dos lados cuenta, así que sumamos las dos colas; como la normal es simétrica, basta calcular una y multiplicarla por dos.

¿Y cómo decidimos con el valor \(p\)? Comparándolo con \(\alpha\). Recordemos que \(\alpha\) es la cantidad de \enquote{rareza} que estamos dispuestos a tolerar antes de desconfiar de \(H_0\). Si el valor \(p\) es más chico que \(\alpha\), nuestra muestra es más rara de lo que estábamos dispuestos a aceptar bajo \(H_0\), y rechazamos. Si es más grande, no rechazamos. 
\begin{gather*}
\text{Si } p<\alpha \Rightarrow \text{rechazamos } H_0,\\
\text{si } p\geq\alpha \Rightarrow \text{no rechazamos.}
\end{gather*}
En nuestro ejemplo, \(p=\num{0,0228}\) es menor que \(\alpha=\num{0,05}\), así que rechazamos \(H_0\): la misma conclusión que antes.

Y no es casualidad que coincida. Los dos métodos son la misma decisión vista de dos formas. Mirémoslo en un solo gráfico: la región de rechazo es todo lo que está a la izquierda de \(z_\alpha=-\num{1,645}\) y encierra un área \(\alpha=\num{0,05}\); el valor \(p\) es el área a la izquierda de nuestro estadístico \(z^*=-2\).

\begin{tikzpicture}
    \begin{axis}[
        width=1.3\linewidth,
        height=5cm,
        axis lines=middle,
        domain=-3.8:3.8,
        samples=150,
        ymin=0,
        enlargelimits=false,
        clip=false,
        ytick=\empty,
        xtick=\empty,
    ]
    \addplot[fill=DiceRed!15, draw=none, domain=-3.8:-1.645] {gauss(x,0,1)} \closedcycle;
    \addplot[fill=DiceRed!45, draw=none, domain=-3.8:-2] {gauss(x,0,1)} \closedcycle;
    \addplot[DiceBlue!80!black, thick] {gauss(x,0,1)};
    \draw[gray] (axis cs:-1.645, 0) -- (axis cs:-1.645, {gauss(-1.645,0,1)});
    \draw[Mahogany, thick] (axis cs:-2, 0) -- (axis cs:-2, {gauss(-2,0,1)});
    \node[below, font=\small] at (axis cs:0, 0) {\(0\)};
    \node[below, font=\small] at (axis cs:-1.645, 0) {\(z_\alpha\)};
    \node[below, Mahogany, font=\small, yshift=-11pt] at (axis cs:-2, 0) {\(z^*\)};
    \node[DiceRed!50!black, font=\small] at (axis cs:-1.8, 0.02) {\(\alpha\)};
    \node[DiceRed!60!black, font=\small] at (axis cs:-2.8, 0.03) {\(p\)};
    \end{axis}
\end{tikzpicture}

Como \(-2\) está más a la izquierda que \(-\num{1,645}\), su cola es más chica que toda la región de rechazo. Y ahí está el truco: el estadístico cae \emph{dentro} de la región de rechazo exactamente cuando el área que deja en la cola (su valor \(p\)) es menor que \(\alpha\). Rechazar porque el estadístico está en la región de rechazo y rechazar porque \(p<\alpha\) son, entonces, la misma cosa.

El valor \(p\) tiene además una ventaja práctica: no depende de haber fijado \(\alpha\) de antemano. Quien lea nuestro informe puede aplicar su propio umbral. Por eso, con un computador a mano, casi siempre se reporta el valor \(p\).

Una advertencia, porque es el error más común: el valor \(p\) \emph{no} es la probabilidad de que \(H_0\) sea verdadera. Es una probabilidad calculada \emph{suponiendo} que \(H_0\) es verdadera, que es muy distinto\sidenote{después del capítulo de Bayes esto debería sonar conocido: \(\prob\giventhat{\text{datos}}{H_0}\) no es \(\prob\giventhat{H_0}{\text{datos}}\)}.

\section{Toma de decisión}
\labsec{toma-decision}

Lo mecánico ya está: tenemos dos formas de comparar el estadístico de prueba —la región de rechazo y el valor \(p\)— y las dos llevan a la misma decisión. Lo que falta es lo más importante y lo que más cuesta: interpretar bien esa decisión.

Cuando rechazamos \(H_0\), estamos diciendo que la muestra es difícil de explicar si \(H_0\) fuera verdadera, tan difícil que preferimos pensar que \(H_0\) es falsa. En el ejemplo de las botellas rechazamos: los datos son convincentes de que se rellena de menos.

Cuando \emph{no} rechazamos, hay que cuidar las palabras. Fíjese que en todo el capítulo dijimos \enquote{no rechazar} y no \enquote{aceptar}\sidenote{por eso la tabla del comienzo decía \enquote{aceptamos} entre comillas mentales}. No es lo mismo. No rechazar \(H_0\) no quiere decir que hayamos probado que es verdadera; quiere decir que la muestra no nos dio evidencia suficiente para descartarla. Ausencia de evidencia no es evidencia de ausencia.

\begin{pausa}
Un test da \enquote{no significativo}: no rechazamos \(H_0\). ¿Probamos que \(H_0\) es verdadera, o sólo que no juntamos evidencia en su contra? ¿Cambia su respuesta si la muestra era de \num{5} datos en vez de \num{5000}?
\end{pausa}

Sirve pensarlo como un juicio. El veredicto \enquote{no culpable} no significa \enquote{probadamente inocente}, sino que la fiscalía no logró probar la culpabilidad más allá de toda duda razonable. La hipótesis nula es el acusado: se la presume cierta hasta que la evidencia —la muestra— demuestre lo contrario, y \(\alpha\) es nuestra \enquote{duda razonable}.

Por eso convivimos siempre con los dos errores del comienzo del capítulo. Rechazar cuando \(H_0\) era cierta es el error tipo I, y su probabilidad la fijamos nosotros: es \(\alpha\). No rechazar cuando \(H_0\) era falsa es el error tipo II, con probabilidad \(\beta\), y es más escurridizo porque depende de cuál sea la verdad. Bajar \(\alpha\) para no rechazar a la ligera sube \(\beta\): nos volvemos más conservadores y dejamos pasar más \(H_0\) falsas. No hay almuerzo gratis.

Al reportar, lo honesto es decir las tres cosas: qué hipótesis contrastamos, con qué significancia, y qué obtuvimos. Un \enquote{rechazamos \(H_0\) al \qty{5}{\percent} de significancia (\(z^*=-2\), \(p=\num{0,0228}\))} dice mucho más que un escueto \enquote{dio significativo}.

Dos advertencias finales para no salir a hacer daño con esto. Primero, \enquote{significativo} no quiere decir \enquote{importante}: con una muestra enorme, una diferencia minúscula y sin ninguna relevancia práctica puede dar un valor \(p\) chiquitito. Que la embotelladora rellene \qty{0,1}{\milli\liter} de menos es detectable con suficientes botellas, pero a nadie le importa. Segundo, si uno prueba muchas hipótesis hasta que una \enquote{dé}, tarde o temprano alguna va a dar por puro azar\sidenote{si hacemos \num{20} pruebas a un \qty{5}{\percent}, esperamos que una salga significativa aunque todas las \(H_0\) sean verdaderas}: la significancia se gana antes de mirar los datos, no después.